\scanpage{9}
Enfin, on a une famille de scindages \inserted{partiels}
\[
(25)\qquad\Psi^{\deleted{\unsure{\mathrm{ext}}}}=(\Psi_j)_{j\in J}:\bigl(\widetilde X'(j)\longrightarrow\operatorname{Hom}(\mathcal G'_j,\mathcal E'_{\deleted{\unsure{X'}}})\bigr)
\]
(où cette fois $\mathcal E'$ désigne l'image inverse
de $\mathcal G'\subset\mathcal G_{X'}$ dans $\mathcal E_{X'}$, \inserted{et $\mathcal G'_j$ le stab. de $j$ dans $\mathcal G'$}), compatible
avec les opérations de $\mathcal E'$ de part et
d'autre, \deleted{\unsure{se prolongeant} dans une famille de
scindages mod op. de $\pi_1$}
\[
\deleted{(26)\qquad\Psi^{\mathrm{ext}}=(\Psi_j^{\mathrm{ext}})_{j\in J}:\bigl(\widetilde X'(J)\longrightarrow\operatorname{Hom}(\mathcal G'_j,\mathcal E')/\operatorname{Int}(\pi_1,\mathcal E')\bigr)}
\]
\deleted{également compatible avec op. de $\mathcal E'$.}

\textcircled{6} \uuline{Récapitulons.} On se donne : \deleted{$X$
schéma de t.f. sur $\overline{\mathbf Q}$, un divi\unsure{-}}
\[
\mathrm I\quad\left\{\begin{array}{l}
X\text{ de t.f. sur }\overline{\mathbf Q}\\
D=\sum_{i\in I}D_i\text{ diviseur effectif sur }X\text{ à comp.}\\
\qquad\text{simples, }X\text{ lisse en les pts minimaux}\\
\qquad\text{de }\operatorname{supp}D\\
J\subset X'(\overline{\mathbf Q}),\quad\text{où }X'=X-D
\end{array}\right.
\]
On pose\margin{\unsure{(}}
\[
\mathrm{II}\quad\left\{\begin{array}{l}
\mathcal G'=\text{gr des autom du schéma \inserted{absolu} }X\\
\qquad\text{qui invarie }D\text{ et }J\\
G'=\text{ss-groupe de }\mathcal G'\text{ formé des }\overline{\mathbf Q}\text{-autom.}\\
\mathbf T(\overline{\mathbf Q})=\varprojlim_n\mu_n(\overline{\mathbf Q})\\[1ex]
\Gamma=\operatorname{Gal}(\overline{\mathbf Q}/\mathbf Q)
\end{array}\right.
\]
\mspage{8}
